\subsection{Related Work}
\label{rasc:sec:related-work}

\noindent \textbf{Push-based Frameworks.}
\abs{} shares semantic roots with ROS actionlib~\cite{ros_actionlib}, gRPC streams~\cite{grpc_streaming}, and MQTT~\cite{mqtt_v5}, which provide goal-feedback-result lifecycles or push-based notifications. However, these frameworks typically require device-side protocol changes. \abs{} is minimal.

\noindent \textbf{Scheduling and Action Dependencies.}  
A range of systems provide abstractions for coordinating IoT and cyber-physical devices. 
Gaia~\cite{gaia}, Bundle~\cite{vicaire2010bundle}, dSpace~\cite{fu2021dspace}, and DepSys~\cite{DepSys} offer middleware or dependency-based programming, but at coarse granularity without action-level progress or dynamic scheduling. 
TransActuations~\cite{Transactuations}, Rivulet~\cite{rivulet}, and IoTRepair~\cite{iotrepair2020} strengthen safety through transactional or rollback semantics, yet operate only at the level of instantaneous or fixed-length actions or routines. 
SafeHome~\cite{SafeHomeEuroSys2021} and Hades~\cite{hades} handle routine conflicts but not fine-grained dependencies. 
In contrast, \sysname{} captures dependencies within as well as across routines, and our schedulers ensure safety and serial equivalence even as actions progress unpredictably.

\noindent \textbf{Dynamic Rescheduling.}  
Classical resource-reclaiming
in real-time systems~\cite{shen1993resource,manimaran1997new,gupta2000new} captures slack from early tasks 
but assumes fixed DAGs and ignores late completions. 
Our adaptation of RV and STF extends these ideas to IoT action graphs with cross-routine serialization constraints, handling both early and delayed completions without violating safety.

\noindent \textbf{Observability and Failure Detection.}  
IoT systems often rely on coarse timeout-based detection or assume instantaneous actions~\cite{kashef2021wireless}, which is unsuitable for long-running tasks. 
SafeHome~\cite{SafeHomeEuroSys2021} and Hades~\cite{hades} provide some exception handling but not progress-aware detection. 
Our adaptive polling provides fine-grained observability of start/complete events, enabling efficient and accurate failure detection without overloading the system.
While pub/sub systems~\cite{gryphon,carzaniga2001design,aguilera1999matching,segall1997elvin} are often proposed as alternatives, they are best suited for settings with many subscribers across diverse topics. \sysname{} runs on a single central hub that subscribes to all topics, so pub/sub offers little benefit, and polling remains the core challenge it addresses.
The histogram bucket problem~\cite{scott1979optimal,freedman1981histogram,knuth2006optimal,scargle2013studies,jagadish1998optimal} is analogous: both tasks require selecting partition ranges (bins vs. polling intervals) under uncertainty, differing only in the optimization criterion—e.g., minimizing error in density estimation versus balancing observability cost and timeliness.

\noindent \textbf{Commercial Systems.}  
Commercial platforms such as Alexa~\cite{Alexa}, Google Home~\cite{GoogleHome}, SmartThings~\cite{SmartThings}, and Home Assistant~\cite{HomeAssistant} support limited forms of automation and do not guarantee conflict freedom or robust failure handling. 
\abs{} generalizes them for expressive DAG-based routines, with progress awareness and principled scheduling guarantees.
